% Options for packages loaded elsewhere
\PassOptionsToPackage{unicode}{hyperref}
\PassOptionsToPackage{hyphens}{url}
\PassOptionsToPackage{dvipsnames,svgnames,x11names}{xcolor}
%
\documentclass[
  10pt,
]{article}
\usepackage{amsmath,amssymb}
\usepackage{iftex}
\ifPDFTeX
  \usepackage[T1]{fontenc}
  \usepackage[utf8]{inputenc}
  \usepackage{textcomp} % provide euro and other symbols
\else % if luatex or xetex
  \usepackage{unicode-math} % this also loads fontspec
  \defaultfontfeatures{Scale=MatchLowercase}
  \defaultfontfeatures[\rmfamily]{Ligatures=TeX,Scale=1}
\fi
\usepackage{lmodern}
\ifPDFTeX\else
  % xetex/luatex font selection
  \setmainfont[]{DejaVu Serif}
  \setsansfont[]{DejaVu Sans}
  \setmonofont[]{DejaVu Sans Mono}
  \setmathfont[]{Latin Modern Math}
\fi
% Use upquote if available, for straight quotes in verbatim environments
\IfFileExists{upquote.sty}{\usepackage{upquote}}{}
\IfFileExists{microtype.sty}{% use microtype if available
  \usepackage[]{microtype}
  \UseMicrotypeSet[protrusion]{basicmath} % disable protrusion for tt fonts
}{}
\makeatletter
\@ifundefined{KOMAClassName}{% if non-KOMA class
  \IfFileExists{parskip.sty}{%
    \usepackage{parskip}
  }{% else
    \setlength{\parindent}{0pt}
    \setlength{\parskip}{6pt plus 2pt minus 1pt}}
}{% if KOMA class
  \KOMAoptions{parskip=half}}
\makeatother
\usepackage{xcolor}
\usepackage[margin=2.0cm]{geometry}
\usepackage{longtable,booktabs,array}
\usepackage{calc} % for calculating minipage widths
% Correct order of tables after \paragraph or \subparagraph
\usepackage{etoolbox}
\makeatletter
\patchcmd\longtable{\par}{\if@noskipsec\mbox{}\fi\par}{}{}
\makeatother
% Allow footnotes in longtable head/foot
\IfFileExists{footnotehyper.sty}{\usepackage{footnotehyper}}{\usepackage{footnote}}
\makesavenoteenv{longtable}
\setlength{\emergencystretch}{3em} % prevent overfull lines
\providecommand{\tightlist}{%
  \setlength{\itemsep}{0pt}\setlength{\parskip}{0pt}}
\setcounter{secnumdepth}{-\maxdimen} % remove section numbering
% doc_header.tex -- LaTeX preamble additions used by build_docs.sh for every
% Markdown -> LaTeX -> PDF conversion (supplements, README, submission files).
\usepackage{etoolbox}
% longtable may not be loaded by pandoc when a document has no table
\usepackage{longtable}
% bold vectors: \bm -> \symbf under unicode-math (XeLaTeX), bm package otherwise
\makeatletter
\@ifpackageloaded{unicode-math}{\providecommand{\bm}[1]{\symbf{#1}}}{\usepackage{bm}}
\makeatother
% smaller type in long tables so that wide audit tables fit the text width;
% inside a table a line break is allowed after an underscore of plain text
% (pandoc writes it as \_), so that long identifiers wrap inside their column
% rows are set a little apart, so that symbols with indices do not touch
\AtBeginEnvironment{longtable}{\footnotesize\renewcommand{\arraystretch}{1.45}%
  \let\sptunderscore\_\def\_{\sptunderscore\allowbreak}}
\setlength{\LTcapwidth}{\textwidth}
% allow line breaks after slashes and hyphens inside table cells
\usepackage{xurl}
% code blocks: break over-long lines instead of running into the margin.
% Plain ``` blocks are typeset by pandoc as {verbatim}, highlighted ones
% (```bash) as {Highlighting}; both are re-declared as breakable fancyvrb
% environments.  The continuation line is marked with a small hooked arrow.
\usepackage{fvextra}
\RecustomVerbatimEnvironment{verbatim}{Verbatim}{breaklines,breakanywhere,fontsize=\small}
\makeatletter
\@ifundefined{Highlighting}{}{%
  \RecustomVerbatimEnvironment{Highlighting}{Verbatim}{breaklines,breakanywhere,fontsize=\small,commandchars=\\\{\}}}
\makeatother
% inline code (pandoc writes it as \texttt{...}): never hyphenate -- an inserted
% hyphen could be read as part of a file name -- but allow a line break after
% an underscore, slash, comma, colon or equals sign, so that long paths and
% column lists wrap inside the text width and inside table columns.
% Inside a formula (e.g. the name of a notebook symbol in an equation) the same
% command must first leave math mode, which \text of amsmath does.
\ExplSyntaxOn
\tl_new:N \l__spt_code_tl
\cs_new_protected:Npn \__spt_texttt:n #1
  {
    \tl_set:Nn \l__spt_code_tl {#1}
    \regex_replace_all:nnN { ([\/\,\:\=]) } { \1 \c{allowbreak} } \l__spt_code_tl
    \regex_replace_all:nnN { \c{_} } { \c{_} \c{allowbreak} } \l__spt_code_tl
    \group_begin: \ttfamily \hyphenchar \font = -1 \scan_stop: \l__spt_code_tl \group_end:
  }
\RenewDocumentCommand \texttt { m }
  {
    \mode_if_math:TF { \text { \__spt_texttt:n {#1} } } { \__spt_texttt:n {#1} }
  }
\ExplSyntaxOff
% A figure is written in the Markdown source as an image paragraph followed by
% a caption paragraph ("**Figure S1.** ..."), so that the caption is also
% visible on GitHub.  An image that starts a paragraph is centered.
\usepackage{graphicx}% pandoc loads it only for documents that contain an image
\NewCommandCopy\sptincludegraphics\includegraphics
\RenewDocumentCommand\includegraphics{O{} m}{%
  \ifvmode\begin{center}\sptincludegraphics[#1]{#2}\end{center}%
  \else\sptincludegraphics[#1]{#2}\fi}
% keep figures where they are referenced
\usepackage{float}
\floatplacement{figure}{H}
% running header with the document title
\usepackage{fancyhdr}
\pagestyle{fancy}
\fancyhf{}
\fancyhead[L]{\small PROJECT spacetimes --- documentation and provenance}
\fancyhead[R]{\small\thepage}
\renewcommand{\headrulewidth}{0.4pt}
% 16x16 matrices: tighter columns so that a full matrix fits the text width
\setlength{\arraycolsep}{3.2pt}
% The documents of this project contain long tagged symbols, file paths and
% identifiers that cannot be hyphenated.  A paragraph is therefore allowed to
% have looser lines before TeX gives up and lets a line stick out of the margin.
\AtBeginDocument{\tolerance=2000 \emergencystretch=5em \hbadness=2000}
\ifLuaTeX
  \usepackage{selnolig}  % disable illegal ligatures
\fi
\IfFileExists{bookmark.sty}{\usepackage{bookmark}}{\usepackage{hyperref}}
\IfFileExists{xurl.sty}{\usepackage{xurl}}{} % add URL line breaks if available
\urlstyle{same}
\hypersetup{
  pdftitle={PROJECT spacetimes --- README},
  pdfauthor={Patrick Lee Nash, Ph.D.},
  colorlinks=true,
  linkcolor={blue},
  filecolor={Maroon},
  citecolor={blue},
  urlcolor={blue},
  pdfcreator={LaTeX via pandoc}}

\title{PROJECT spacetimes --- README}
\usepackage{etoolbox}
\makeatletter
\providecommand{\subtitle}[1]{% add subtitle to \maketitle
  \apptocmd{\@title}{\par {\large #1 \par}}{}{}
}
\makeatother
\subtitle{A singlet superfluid of a 16-component Dirac field on 4+4 and
on 7+1 Minkowski spacetime: where everything is, and how to check it}
\author{Patrick Lee Nash, Ph.D.}
\date{November 12, 2026}

\begin{document}
\maketitle

\hypertarget{what-this-project-is}{%
\section{1. What this project is}\label{what-this-project-is}}

The project asks one question for two flat eight-dimensional spacetimes,
one with four space-like and four time-like dimensions (4+4) and one
with seven space-like dimensions and one time-like dimension (7+1). Can
a 16-component Dirac field form a spin-singlet superfluid whose pairing
is mediated by a new scalar boson, the \emph{Vorton-Scalar}? And can
that superfluid be the Higgs field, the dark matter, or the remedy of
the Hubble tension?

The coordinates are \(x0,\dots,x7\) in both spacetimes: \(x0\) is a
hidden space-like coordinate, \(x1,x2,x3\) are ordinary space, \(x4\) is
ordinary time, and \(x5,x6,x7\) are additional coordinates (time-like in
4+4, space-like in 7+1).

\hypertarget{the-answers-in-brief}{%
\section{2. The answers in brief}\label{the-answers-in-brief}}

The complete statements, with their proofs, their numbers and their
limits, are in the two manuscripts; each section of the manuscripts ends
with a subsection ``What has been established, and what has not''.

\begin{itemize}
\tightlist
\item
  \textbf{The matrices.} Each spacetime has eight \emph{real}
  \(16\times16\) Dirac matrices. Their anticommutation relations, and
  the fact that their scaled commutators generate the spin groups (with
  the matrices themselves: \(\mathrm{Pin}(4,4)\) and
  \(\mathrm{Pin}(7,1)\)), are proved and verified by 106 machine checks
  and by an independent second program.
\item
  \textbf{The superfluid.} Established within stated approximations: the
  field has a pair of spin 0 that is a scalar of ordinary spacetime; a
  complex scalar field, the Vorton-Scalar, binds this pair, and its
  neutral current selects it; at any finite density the paired state
  lies below the unpaired one. The singlet along the hidden coordinate
  is a ground state, but one of a degenerate family.
\item
  \textbf{The Higgs hypothesis.} Not established. On 4+4 it has no
  geometric basis. On 7+1 one structural condition is met (the order
  parameter is a multiplet of the rotations of the four extra axes, and
  its vacuum value breaks \(SO(4)\) to \(SO(3)\), the pattern of the
  Higgs vacuum), but the dynamical conditions are not met where the
  model can be solved. In both spacetimes a condensate in empty space
  was found only at couplings at which the methods used are a guide and
  not a proof.
\item
  \textbf{The Hubble tension.} Not explained: simulations with a
  purpose-built Einstein-Boltzmann solver show that the dark radiation
  of the model raises the Hubble constant only at a price that the data
  do not allow.
\item
  \textbf{Dark matter.} A superfluid that could act on galactic scales
  erases the observed small-scale structure. Heavy bound pairs as cold
  dark matter are neither supported nor excluded, and make no use of the
  superfluid.
\end{itemize}

\hypertarget{where-to-start-reading}{%
\section{3. Where to start reading}\label{where-to-start-reading}}

\begin{longtable}[]{@{}
  >{\raggedright\arraybackslash}p{(\columnwidth - 2\tabcolsep) * \real{0.4239}}
  >{\raggedright\arraybackslash}p{(\columnwidth - 2\tabcolsep) * \real{0.5761}}@{}}
\toprule\noalign{}
\begin{minipage}[b]{\linewidth}\raggedright
If you want
\end{minipage} & \begin{minipage}[b]{\linewidth}\raggedright
Read
\end{minipage} \\
\midrule\noalign{}
\endhead
\bottomrule\noalign{}
\endlastfoot
the whole argument for one spacetime &
\texttt{manuscript/16\ Component\ Dirac\ Fermion\ Condensate\ Singlet\ Superfluid\ on\ 4+4\ Spacetime.pdf},
or \texttt{...\ on\ 7+1\ Spacetime.pdf} \\
the proofs about the matrices & \texttt{dirac\ matrices.pdf} \\
the Vorton-Scalar and the superfluid in detail &
\texttt{supplement/S1\_vorton\_scalar\_and\_superfluid.pdf} \\
the Higgs hypothesis & \texttt{supplement/S2\_higgs\_hypothesis.pdf} \\
the cosmological simulations & \texttt{supplement/S3\_cosmology.pdf} \\
the effective action of Cornwall, Jackiw and Tomboulis &
\texttt{supplement/S4\_cjt\_effective\_action.pdf} \\
what every symbol means, and how it differs from the symbols of the
original files & \texttt{provenance/P1\_master\_symbols.pdf} \\
how to compile and run the code & \texttt{provenance/P8\_code.pdf} \\
\end{longtable}

Every document exists in three versions with the same name: Markdown
(\texttt{.md}), LaTeX (\texttt{.tex}) and PDF (\texttt{.pdf}).

\hypertarget{what-is-in-the-repository}{%
\section{4. What is in the repository}\label{what-is-in-the-repository}}

\begin{longtable}[]{@{}
  >{\raggedright\arraybackslash}p{(\columnwidth - 2\tabcolsep) * \real{0.3068}}
  >{\raggedright\arraybackslash}p{(\columnwidth - 2\tabcolsep) * \real{0.6932}}@{}}
\toprule\noalign{}
\begin{minipage}[b]{\linewidth}\raggedright
Path
\end{minipage} & \begin{minipage}[b]{\linewidth}\raggedright
Content
\end{minipage} \\
\midrule\noalign{}
\endhead
\bottomrule\noalign{}
\endlastfoot
\texttt{manuscript/} & the two manuscripts (REVTeX, PDF, Markdown) and
their bibliography \\
\texttt{dirac\ matrices.*} & provenance ``dirac matrices'': definitions,
theorems, proofs and worked examples for both sets of matrices \\
\texttt{supplement/} & Supplements S1 to S4 \\
\texttt{provenance/} & P0 (how commits reach GitHub), P1 (MASTER
provenance of the symbols), P2 (the PDF files used), P3 (the notebooks
used), P8 (the code), P9 (the error log); \texttt{symbols.csv};
\texttt{inputs\_sha256.txt}; \texttt{generated/} (fragments made by the
programs) \\
\texttt{code/spacetimes-rs/} & the Rust program that computes everything
(no external libraries; 63 unit tests) \\
\texttt{code/scripts/} & Python scripts: independent verification,
tables, figures, building of the documents, checks of the symbols \\
\texttt{tools/} & reading the Mathematica notebooks without
Mathematica \\
\texttt{data/} & every data file that a program wrote, each with a
header that states the program, the command, the content and the
units \\
\texttt{figures/} & every figure, as PDF and PNG \\
\texttt{docs\_src/} & the sources of the documents, with placeholders
for the generated numbers \\
\end{longtable}

Not in the repository: the author's Mathematica notebooks, the nine
published papers of the project folder, and the task specification. The
provenances identify these files by their checksums
(\texttt{provenance/inputs\_sha256.txt}).

\hypertarget{how-to-check-the-results}{%
\section{5. How to check the results}\label{how-to-check-the-results}}

The short form (details and explanations: \texttt{provenance/P8\_code}):

\begin{verbatim}
cd "/Users/nsh/Developer/github/spacetimes/code/spacetimes-rs"
cargo test --release                          # 63 unit tests
cd "/Users/nsh/Developer/github/spacetimes"
code/spacetimes-rs/target/release/spacetimes matrices           # the 106 machine checks
python3 code/scripts/verify_matrices_independent.py             # the independent program
python3 code/scripts/check_symbols.py                           # the checks of the symbols
\end{verbatim}

The first command needs the Rust compiler; \texttt{provenance/P8\_code}
says how to install it and how to recompute every data file, table and
figure.

\hypertarget{rules-that-the-project-follows}{%
\section{6. Rules that the project
follows}\label{rules-that-the-project-follows}}

\begin{itemize}
\tightlist
\item
  Every object has one symbol; symbols of the 4+4 spacetime carry the
  tag 1644, those of the 7+1 spacetime the tag 1671, and those of the
  Vorton-Scalar the suffix VS (\texttt{provenance/P1\_master\_symbols}).
\item
  No number in a document is typed by hand: programs write data files,
  and scripts insert the numbers into the documents.
\item
  Every error that was found is recorded, with the way it was found and
  its correction (\texttt{provenance/P9\_error\_log}).
\item
  What is not established is said, in every section.
\end{itemize}

\hypertarget{licence}{%
\section{7. Licence}\label{licence}}

The source code is distributed under the GNU General Public License,
version 3 or later (file \texttt{LICENSE}).

\end{document}
